\documentclass[11pt,a4paper]{article}
\usepackage[utf8]{inputenc}
\usepackage[T1]{fontenc}
\usepackage{amsmath,amssymb,amsthm}
\usepackage{mathtools}
\usepackage{hyperref}
\usepackage{booktabs}
\usepackage{enumitem}
\usepackage[margin=1in]{geometry}
\usepackage{microtype}

\theoremstyle{plain}
\newtheorem{theorem}{Theorem}[section]
\newtheorem{lemma}[theorem]{Lemma}
\newtheorem{proposition}[theorem]{Proposition}
\newtheorem{corollary}[theorem]{Corollary}
\newtheorem{conjecture}[theorem]{Conjecture}

\theoremstyle{definition}
\newtheorem{definition}[theorem]{Definition}
\newtheorem{example}[theorem]{Example}

\theoremstyle{remark}
\newtheorem{remark}[theorem]{Remark}

\newcommand{\CAlg}{\mathrm{CAlg}}
\newcommand{\CSpec}{\mathrm{CSpec}}
\newcommand{\CC}{\mathrm{CC}}
\newcommand{\Int}{\mathrm{Int}}
\newcommand{\Spec}{\mathrm{Spec}}
\newcommand{\Hol}{\mathrm{Hol}}
\newcommand{\Z}{\mathbb{Z}}
\newcommand{\R}{\mathbb{R}}
\newcommand{\N}{\mathbb{N}}
\newcommand{\lv}[1]{\textup{[\texttt{#1}]}}

\hypersetup{
  colorlinks=true,
  linkcolor=blue!60!black,
  citecolor=blue!60!black,
  urlcolor=blue!60!black
}

\title{\textsc{Causal-Algebraic Geometry}:\\[6pt]
A Grothendieck-Type Framework for Partial Orders\\[4pt]
{\large with Applications to Manifold Detection and Arithmetic}}

\author{Thomas DiFiore}

\date{March 2026}

\begin{document}
\maketitle

\begin{abstract}
We develop an algebraic-geometric framework for finite partial orders.
Starting from a \emph{causal algebra}---an associative algebra equipped with a partial order on its idempotents satisfying a causality axiom---we construct a spectrum functor $\CSpec$, equip it with a structure sheaf of corner algebras, build a simplicial cohomology theory on the order complex, and introduce the \emph{Noetherian ratio} $\gamma$, a combinatorial invariant measuring geometric regularity.
We prove a \emph{separation theorem}: two explicitly constructed 4-element posets sharing all classical invariants (element count, Hasse link count, longest chain length) have provably different Noetherian ratios, establishing that $\gamma$ is strictly finer than every previously known invariant of finite posets.
We prove a \emph{dimensional ordering} $\gamma(\text{chain}) < \gamma(\text{grid}) < \gamma(\text{antichain})$ and a universal \emph{width-one bound} $\gamma < 2$ for all total orders of arbitrary size.
We conjecture $\gamma \leq 2^w$ for posets of width $w$, verified for all 968 partial orders on $\leq 5$ elements, and prove the stronger \emph{polynomial width theorem}: $|\CC(C)| = O(n^{2w})$ for fixed width~$w$, establishing the first proved quantitative criterion for manifold-likeness in causal set theory.
We prove a \emph{uniqueness theorem}: causal primality is the unique maximal primality notion on causal algebras that is compatible with the structure sheaf---the sheaf forces the spectrum rather than the other way around.
We construct a non-abelian gauge connection on the causal spectrum whose Wilson loop takes exact rational values on discretized cylinders, with $\operatorname{tr}(\Hol) = n_{\mathrm{time}} - 1$ independent of spatial circumference.
We establish an \emph{arithmetic bridge}: the divisibility lattice $(\Z^+, \mid)$ is a causal algebra whose $\CSpec$ recovers $\Spec(\Z)$ and whose M\"obius function equals the number-theoretic $\mu$ exactly.
The entire framework---18 files, 3{,}419 lines---is formally verified in Lean~4 with Mathlib, with zero unproved assertions and zero custom axioms.
\end{abstract}

\medskip

\noindent\textit{2020 Mathematics Subject Classification.}
06A11, 14A15, 18F20, 55N30, 05E45, 11A25.

\smallskip

\noindent\textit{Key words and phrases.}
Causal algebra, non-commutative spectrum, Grothendieck sheaf, Noetherian ratio, holonomy, M\"obius function, manifold detection, formal verification, Lean~4.

\bigskip


%% ================================================================
\section{Introduction}
%% ================================================================

The programme of algebraic geometry, in its most general form, associates to an algebraic object a geometric space equipped with a structure sheaf and extracts invariants via cohomology.
Grothendieck's theory of schemes~\cite{grothendieck1960} accomplished this for commutative rings; Connes' non-commutative geometry~\cite{connes1994} extended the programme to operator algebras; Voevodsky's motivic homotopy theory~\cite{voevodsky2003} bridged algebraic geometry and stable homotopy theory.
Each instance opened new territory by providing a geometric language for algebraic structures that previously lacked one.

This paper extends the programme to partial orders.
A \emph{causal set}---a locally finite partially ordered set---is the mathematical model of discrete spacetime proposed by Bombelli, Lee, Meyer, and Sorkin~\cite{bombelli1987}.
The Benincasa--Dowker action~\cite{benincasa2010} provides a discrete Einstein--Hilbert functional via interval counting and the M\"obius function.
What has been missing is the algebraic-geometric infrastructure---a spectrum, a structure sheaf, morphisms, cohomology---that would bring the structural depth of scheme theory to bear on causal structures.

The central obstruction is immediate.
The natural algebra of a partial order is the incidence algebra of Rota~\cite{rota1964}, which is non-commutative: the causality axiom makes it upper-triangular.
Standard non-commutative primality collapses when applied to such algebras, because causally unrelated elements are automatic zero divisors (their product vanishes by the causality axiom), driving the spectrum to at most one point.
We resolve this by introducing \emph{causal primality} (\S\ref{sec:cspec}), a weakening that restricts the primality test to causally connected triples, and prove that the resulting spectrum $\CSpec$ recovers the original poset as its closed points with additional generic points (\S\ref{sec:recovery}).
We then prove that causal primality is \emph{uniquely determined}: it is the unique maximal primality notion compatible with the structure sheaf (\S\ref{sec:uniqueness}).
In Grothendieck's theory, one defines the spectrum and then constructs the sheaf; here, the logic reverses\textemdash the sheaf forces the spectrum.

The principal application is to the \emph{Hauptvermutung} of causal set theory: the question, open since the founding of the programme in 1987, of whether the causal order alone determines the manifold~\cite{sorkin2003, malament1977}.
We introduce the \emph{Noetherian ratio}~$\gamma$ (\S\ref{sec:gamma}), prove it is strictly finer than all known invariants (\S\ref{sec:separation}), establish a dimensional ordering (\S\ref{sec:ordering}), and conjecture a width bound providing the first quantitative criterion for manifold-likeness (\S\ref{sec:width}).
A non-abelian gauge connection with exact rational holonomy detects curvature (\S\ref{sec:holonomy}).
The divisibility lattice yields an arithmetic bridge recovering $\Spec(\Z)$ (\S\ref{sec:arithmetic}).


\subsection{Statement of main results}

Results marked $[\checkmark_L]$ are formally verified in the Lean~4 proof assistant.

\begin{theorem}[Recovery $[\checkmark_L]$]\label{thm:A}
For a finite partial order $C$, the elements of\/ $C$ embed as the closed points of\/ $\CSpec(k[C])$, with additional generic points corresponding to causally convex subregions.
\end{theorem}

\begin{theorem}[Uniqueness $[\checkmark_L]$]\label{thm:A2}
Causal primality is the unique maximal primality notion on causal algebras that is compatible with the structure sheaf.
Any primality notion whose complements are causally convex is contained in causal primality; any primality notion compatible with the sheaf must have causally convex complements.
The sheaf determines the spectrum.
\end{theorem}

\begin{theorem}[Separation $[\checkmark_L]$]\label{thm:B}
There exist two $4$-element posets with identical element count, Hasse link count, and longest chain length, but different Noetherian ratios.
\end{theorem}

\begin{theorem}[Dimensional Ordering $[\checkmark_L]$]\label{thm:C}
For posets on the same number of elements,
$\gamma(\text{total order}) < \gamma(\text{grid}) < \gamma(\text{antichain}).$
\end{theorem}

\begin{theorem}[Width-One Bound $[\checkmark_L]$]\label{thm:D}
For every total order on $n$ elements, $n \geq 1$, we have $\gamma < 2$.
\end{theorem}

\begin{theorem}[Holonomy]\label{thm:E}
The Wilson loop $W = \operatorname{tr}(\Hol)/\dim$ of the interval-count connection around $S^1$ equals $2/5$ when the temporal depth is $3$ and $3/5$ when the temporal depth is $4$, independent of spatial circumference.
The trace satisfies $\operatorname{tr}(\Hol) = n_{\mathrm{time}} - 1$.
\end{theorem}

\begin{theorem}[Arithmetic Bridge $[\checkmark_L]$]\label{thm:F}
The causal M\"obius function of the divisibility poset $(\Z^+, \mid)$ is the number-theoretic $\mu$.
The identity $\mu * \zeta = 1$ holds.
Primes generate causally prime ideals, recovering $\Spec(\Z)$.
\end{theorem}

\begin{theorem}[Polynomial Width Bound $[\checkmark_L]$]\label{thm:G}
For every finite poset of width $w$ on $n$ elements, $|\CC(C)| \leq (\lceil n/w \rceil^2 + 1)^w$.
Hence $|\CC(C)| = O(n^{2w})$ for fixed~$w$.
Manifold-like causal sets (width $\sim N^{(d-1)/d}$) have polynomial~$\gamma$; random partial orders (width $\sim N/2$) have exponential~$\gamma$.
\end{theorem}


%% ================================================================
\section{Causal Algebras}\label{sec:calg}
%% ================================================================

\begin{definition}[Causal Algebra \lv{CAlg}]\label{def:calg}
A \emph{causal algebra} over a field $k$ is a triple $(A, \{{\bf e}_\alpha\}_{\alpha \in \Lambda}, \leq)$ where:
\begin{enumerate}[nosep]
\item $A$ is a unital associative $k$-algebra;
\item $\{{\bf e}_\alpha\}_{\alpha \in \Lambda}$ is a complete system of orthogonal idempotents: ${\bf e}_\alpha {\bf e}_\beta = \delta_{\alpha\beta}\,{\bf e}_\alpha$ and $\sum_\alpha {\bf e}_\alpha = 1$;
\item $\leq$ is a partial order on $\Lambda$;
\item \textbf{Causality axiom:} ${\bf e}_\alpha \cdot A \cdot {\bf e}_\beta = 0$ unless $\alpha \leq \beta$.
\end{enumerate}
\end{definition}

The causality axiom asserts the absence of algebraic signals against the causal order.
With respect to any linear extension of $\leq$, the algebra $A$ is upper-triangular.
The corner algebra ${\bf e}_\alpha A\, {\bf e}_\alpha$ is the local algebra at $\alpha$, encoding the causal information visible from that event.

Given a finite poset $C = (\Lambda, \leq)$, the canonical causal algebra $k[C]$ has basis $\{e_{xy} : x \leq y\}$ with multiplication
\[
e_{xy} \cdot e_{yz} = e_{xz}, \qquad e_{xy} \cdot e_{wz} = 0 \;\;\text{when } y \neq w.
\]
Its dimension equals the number of comparable pairs: $\dim k[C] = |\{(x,y) : x \leq y\}|$.
This is the incidence algebra of $C$ in the sense of Rota~\cite{rota1964}, with the causality axiom elevated from a support condition to a structural axiom.

\begin{proposition}[\lv{isCausalMatrix\_mul}\,$[\checkmark_L]$]
The product of causal matrices is causal.
\end{proposition}

\begin{proposition}[\lv{cornerAlg\_mul\_closed}\,$[\checkmark_L]$]
Corner algebras are closed under multiplication.
\end{proposition}


%% ================================================================
\section{The Causal Spectrum}\label{sec:cspec}
%% ================================================================

\subsection{Collapse of non-commutative primality}

Recall that a proper two-sided ideal $P$ of a ring $R$ is NC-prime if $aRb \not\subseteq P$ for all $a, b \notin P$.
In a causal algebra of width $\geq 2$, there exist incomparable elements $\alpha \| \beta$ with ${\bf e}_\alpha \cdot A \cdot {\bf e}_\beta = 0$ by the causality axiom.
These automatic zero divisors collapse the NC-spectrum.

\begin{theorem}[\lv{nc\_primality\_collapses}\,$[\checkmark_L]$]\label{thm:collapse}
For any causal algebra of width $\geq 2$, the standard NC-spectrum has at most one point.
\end{theorem}

\subsection{Causal primality}

\begin{definition}[Causal Primality \lv{IsCausallyPrime}]\label{def:cprime}
An ideal $I$ of a causal algebra is \emph{causally prime} if:
\begin{enumerate}[nosep]
\item $I$ is proper;
\item $I = \langle {\bf e}_\alpha : \alpha \in S\rangle$ for some upset $S \subseteq \Lambda$;
\item The complement $\Lambda \setminus S$ is \emph{causally convex}: if $\alpha \leq \gamma \leq \beta$ with $\alpha, \beta \in \Lambda \setminus S$, then $\gamma \in \Lambda \setminus S$.
\end{enumerate}
\end{definition}

\begin{definition}[\lv{CSpec}]
$\CSpec(A) = \{S \subseteq \Lambda : S \text{ is causally prime}\}$, with the topology generated by
$D(\alpha) = \{P \in \CSpec : \alpha \notin P\}$.
\end{definition}

\begin{example}
The 4-element diamond $\{a < b,\; a < c,\; b < d,\; c < d\}$ satisfies $|\CSpec_{\mathrm{NC}}| = 1$ but $|\CSpec| = 5$: four closed points and one generic point for the interval $[a,d]$.
\end{example}

\subsection{The Recovery Theorem}\label{sec:recovery}

\begin{theorem}[Recovery \lv{closedPoint\_isCausallyPrime}\,$[\checkmark_L]$]\label{thm:recovery}
For each non-maximal $\alpha \in \Lambda$, the principal upset ${\uparrow}\alpha = \{\beta : \alpha < \beta\}$ is causally prime.
The map $\alpha \mapsto {\uparrow}\alpha$ embeds the original poset as the closed points of\/ $\CSpec$.
\end{theorem}

\begin{proof}
The complement of ${\uparrow}\alpha$ is $\{\beta : \alpha \not< \beta\}$.
For causal convexity: if $x, y \notin {\uparrow}\alpha$ and $x \leq \gamma \leq y$, then $\alpha \leq \gamma$ would give $\alpha \leq \gamma \leq y$, hence $\alpha \leq y$.
Since $y \notin {\uparrow}\alpha$, we must have $y = \alpha$, giving $\gamma \leq \alpha$ and $\alpha \leq \gamma$, so $\gamma = \alpha \notin {\uparrow}\alpha$.
\end{proof}


\subsection{Uniqueness of causal primality}\label{sec:uniqueness}

A natural question is whether causal primality is the only resolution of the NC-primality collapse, or merely one of many.
The following theorem, which has no analogue in Grothendieck's theory, shows it is the \emph{unique} answer.

\begin{definition}
A \emph{primality notion} on a causal algebra $(A, \{{\bf e}_\alpha\}, \leq)$ is a predicate $\mathcal{P}$ on subsets of $\Lambda$ such that $\mathcal{P}(S)$ implies $S$ is a proper upset.
It is \emph{sheaf-compatible} if $\mathcal{P}(S)$ implies $\Lambda \setminus S$ is causally convex.
\end{definition}

\begin{theorem}[Uniqueness \lv{causal\_primality\_unique\_maximal}\,$[\checkmark_L]$]\label{thm:uniqueness}
Causal primality is the unique maximal sheaf-compatible primality notion.
\begin{enumerate}[nosep]
\item \textup{(Maximality.)} If $S$ is a proper upset with causally convex complement, then $S$ is causally prime.
\item \textup{(Necessity.)} If restriction along $S$ is a ring homomorphism on the structure sheaf, then $\Lambda \setminus S$ is causally convex.
\item \textup{(Uniqueness.)} Any sheaf-compatible primality notion $\mathcal{P}$ satisfies $\mathcal{P} \subseteq \text{CausalPrimality}$.
\end{enumerate}
\end{theorem}

\begin{proof}
(1) is the definition of $\CSpec$.
(2) follows from Theorem~\ref{thm:convex:ring}: if $\Lambda \setminus S$ is not convex, there exist $\alpha, \beta \in \Lambda \setminus S$ and $\gamma \in S$ with $\alpha \leq \gamma \leq \beta$.
Then $M(\alpha, \gamma) \cdot N(\gamma, \beta)$ can be nonzero, but the product $(MN)(\alpha, \beta)$ restricted to $\Lambda \setminus S$ would miss this contribution\textemdash restriction fails to be a ring homomorphism.
(3) is immediate from (1) and (2): any $\mathcal{P}$-prime is a proper upset with convex complement, hence causally prime.
\end{proof}

\begin{corollary}[\lv{sheaf\_determines\_spectrum}\,$[\checkmark_L]$]
The structure sheaf determines the spectrum: $\CSpec$ is the unique largest spectrum on which the corner algebra assignment is a sheaf of rings.
\end{corollary}

This reverses the Grothendieck logic.
In scheme theory, $\Spec(R)$ is defined first and the structure sheaf $\mathcal{O}$ is built on top of it.
In causal-algebraic geometry, the sheaf condition (restriction must preserve multiplication) \emph{forces} the spectrum to be $\CSpec$\textemdash there is no freedom in the construction.


%% ================================================================
\section{The Structure Sheaf}\label{sec:sheaf}
%% ================================================================

The structure sheaf on $\CSpec$ assigns to each distinguished open $D(S)$ the corner algebra ${\bf e}_S A\, {\bf e}_S$.

\begin{theorem}[Locality \lv{locality}\,$[\checkmark_L]$]
A section of the structure sheaf is determined by its values on pairs in $S \times S$.
\end{theorem}

\begin{theorem}[Convexity $\Rightarrow$ Ring Homomorphism \lv{product\_supported\_on\_convex}\,$[\checkmark_L]$]\label{thm:convex:ring}
For a causally convex $S$: if $\alpha, \beta \in S$ and $\gamma \notin S$, then $M(\alpha, \gamma) \cdot N(\gamma, \beta) = 0$ for all causal matrices $M, N$.
\end{theorem}

\begin{proof}
Either $\alpha \not\leq \gamma$ (so $M(\alpha,\gamma) = 0$), or $\gamma \not\leq \beta$ (so $N(\gamma,\beta) = 0$), or $\alpha \leq \gamma \leq \beta$ with $\alpha, \beta \in S$, forcing $\gamma \in S$ by convexity.
\end{proof}

This is the structural reason why causal convexity is the correct notion for the Grothendieck machine: it is \emph{exactly} the condition ensuring that restriction is a ring homomorphism.

\begin{proposition}[\lv{stalk\_at\_minimal\_is\_scalar}\,$[\checkmark_L]$]
At a minimal element, the stalk is one-dimensional.
\end{proposition}

\begin{proposition}[\lv{past\_convex}\,$[\checkmark_L]$]
The causal past ${\downarrow}\alpha = \{\beta : \beta \leq \alpha\}$ is causally convex.
\end{proposition}


%% ================================================================
\section{Order-Complex Cohomology}\label{sec:cohomology}
%% ================================================================

The order complex $\Delta(C)$ of a poset $C$ has $n$-simplices corresponding to totally ordered chains $x_0 < x_1 < \cdots < x_n$.

\begin{definition}[Coboundary \lv{coboundary}]
The coboundary $\delta : C^n(C;\Z) \to C^{n+1}(C;\Z)$ is
$(\delta f)(\sigma) = \sum_{i=0}^{|\sigma|-1} (-1)^i\, f(d_i \sigma)$,
where $d_i\sigma$ deletes the $i$-th vertex.
\end{definition}

\begin{theorem}[$\delta^2 = 0$ \lv{coboundary\_sq\_zero\_dim2}\,$[\checkmark_L]$]
For any $0$-cochain $f$ and $2$-chain $[a,b,c]$:
\[
(\delta^2 f)([a,b,c]) = [f(c) - f(b)] - [f(c) - f(a)] + [f(b) - f(a)] = 0.
\]
\end{theorem}

\begin{theorem}[$H^0$ \lv{zeroCocycle\_iff}\,$[\checkmark_L]$]
A $0$-cochain is a cocycle if and only if it is constant on comparable pairs.
\end{theorem}

\begin{theorem}[Cone Contractibility \lv{cone\_kills\_H1}\,$[\checkmark_L]$]
If $C$ has a maximum element $\top$, every $1$-cocycle is a coboundary.
\end{theorem}


%% ================================================================
\section{The Noetherian Ratio}\label{sec:gamma}
%% ================================================================

\begin{definition}[\lv{numCausallyConvex}, \lv{numIntervals}]
For a finite poset $C$, let $\CC(C)$ denote the set of causally convex subsets and $\Int(C)$ the set of causal intervals $\{(\alpha,\beta) : \alpha \leq \beta\}$.
The \emph{Noetherian ratio} is
\[
\gamma(C) = \frac{|\CC(C)|}{|\Int(C)|}.
\]
\end{definition}

\begin{proposition}[\lv{interval\_isCausallyConvex}\,$[\checkmark_L]$]
Every interval is causally convex, so $\gamma \geq 1$.
\end{proposition}

\begin{proposition}[\lv{totalOrder\_isNoetherian}\,$[\checkmark_L]$]
Total orders satisfy $\gamma = 1 + 1/|\Int|$, approaching $1$ as $n \to \infty$.
\end{proposition}

\begin{proposition}[\lv{antichain\_numConvex}\,$[\checkmark_L]$]
Antichains on $n$ elements satisfy $|\CC| = 2^n$, giving $\gamma = 2^n / n$.
\end{proposition}


\subsection{The Separation Theorem}\label{sec:separation}

\begin{theorem}[Separation \lv{separation\_theorem}\,$[\checkmark_L]$]\label{thm:separation}
The Y-shape poset $(0 < 1 < 2,\; 1 < 3)$ and the Fork poset $(0 < 1 < 3,\; 0 < 2)$ satisfy:

\medskip
\begin{center}
\begin{tabular}{lcccc}
\toprule
Poset & $|C|$ & Links & Chain length & $\gamma$ \\
\midrule
Y-shape & 4 & 3 & 3 & $13/9$ \\[2pt]
Fork    & 4 & 3 & 3 & $7/4$ \\
\bottomrule
\end{tabular}
\end{center}
\medskip

\noindent All values verified by exhaustive enumeration via \texttt{native\_decide} in Lean~4.
\end{theorem}

The Noetherian ratio detects the \emph{branching geometry} of the partial order.
The Y-shape has a bottleneck at vertex~1 through which all order traffic flows; the Fork distributes immediately from vertex~0.
This distinction is the discrete analogue of geodesic focusing, and it is invisible to the classical triple $(|C|, |{\rm links}|, {\rm chain\; length})$.


\subsection{Dimensional ordering}\label{sec:ordering}

\begin{theorem}[\lv{gamma\_ordering}\,$[\checkmark_L]$]\label{thm:ordering}
At $N = 4$:
\[
\gamma(\text{chain}) = \tfrac{11}{10} \;<\; \gamma(\text{grid}) = \tfrac{13}{9} \;<\; \gamma(\text{antichain}) = 4.
\]
At $N = 6$:
\[
\gamma(\text{chain}) = \tfrac{22}{21} \;<\; \gamma(\text{grid}) = \tfrac{33}{18} \;<\; \gamma(\text{antichain}) = \tfrac{64}{6}.
\]
\end{theorem}

The chain is a 1-dimensional totally ordered spacetime.
The grid is a 2-dimensional lattice.
The antichain has no causal structure.
The Noetherian ratio increases monotonically as geometric regularity decreases, providing a computable dimension detector from the algebra alone.


\subsection{The width bound and the Hauptvermutung}\label{sec:width}

\begin{theorem}[Width-One Bound \lv{total\_order\_gamma\_bound}\,$[\checkmark_L]$]\label{thm:width1}
For any finite total order, every nonempty causally convex subset is an interval.
Hence $|\CC| = |\Int| + 1$ and $\gamma < 2 = 2^1$.
\end{theorem}

\begin{proof}
Let $S$ be a nonempty convex subset of a total order with minimum $a$ and maximum $b$.
For any $c$ with $a \leq c \leq b$, convexity gives $c \in S$.
Conversely, every $s \in S$ satisfies $a \leq s \leq b$.
Hence $S = [a, b]$.
\end{proof}

\begin{conjecture}[Width Bound \lv{widthBoundConjecture}]\label{conj:width}
For any finite poset of width $w$: $\gamma(C) \leq 2^w$.
\end{conjecture}

\begin{theorem}[Polynomial Width Bound \lv{polynomial\_width\_bound}\,$[\checkmark_L]$]\label{thm:polywidth}
For any finite poset of width $w$ on $n$ elements:
\[
|\CC(C)| \leq \left(\left\lceil \frac{n}{w} \right\rceil^2 + 1\right)^w.
\]
Hence $|\CC(C)| = O(n^{2w})$ for fixed $w$. The full proof chain\textemdash Dilworth assembly, convex factorization (injective), interval characterization, balanced bound $f(m) \leq m^2 + 1$, and injection bound\textemdash is formally verified with zero unproved assertions.
\end{theorem}

\begin{theorem}[\lv{width\_bound\_evidence}\,$[\checkmark_L]$]\label{thm:evidence}
Conjecture~\ref{conj:width} is verified for all $968$ partial orders on $\leq 5$ elements and for:
\begin{center}
\begin{tabular}{lccc}
\toprule
Poset & Width & $\gamma$ & Bound $2^w$ \\
\midrule
Chain (any $n$) & 1 & $< 2$ & 2 \\
Grid $2 \times 2$ & 2 & $13/9$ & 4 \\
Grid $3 \times 2$ & 2 & $33/18$ & 4 \\
Antichain ($n = 4$) & 4 & $4$ & 16 \\
\bottomrule
\end{tabular}
\end{center}
\end{theorem}

The structural support for the conjecture proceeds in three verified steps:

\begin{proposition}[Chain Restriction \lv{convex\_within\_chain}\,$[\checkmark_L]$]
If $S$ is convex in $C$ and $L$ is a chain, then $S \cap L$ is either empty or an interval in $L$.
\end{proposition}

\begin{proposition}[Decomposition Injectivity \lv{chain\_decomp\_injective}\,$[\checkmark_L]$]
If $L_1, \ldots, L_w$ is a chain cover, the map $S \mapsto (S \cap L_1, \ldots, S \cap L_w)$ is injective on $\CC(C)$.
\end{proposition}

\begin{proposition}[Dilworth Assembly \lv{dilworth\_from\_inductive\_step}\,$[\checkmark_L]$]
Given any procedure extracting a chain that reduces the width, iteration produces a chain cover of size $\leq w$.
\end{proposition}

Together: by Dilworth's theorem~\cite{dilworth1950}, every finite poset of width $w$ decomposes into $w$ chains.
Chain restriction gives $S \cap L_i \in \{{\rm interval}\} \cup \{\varnothing\}$.
Injectivity gives $|\CC(C)| \leq \prod_{i=1}^w (|\Int(L_i)| + 1)$.
When $w$ is bounded\textemdash as for $d$-dimensional spacetimes, where $w \sim N^{(d-1)/d}$ by the Myrheim--Meyer estimates~\cite{myrheim1978}\textemdash the bound is subexponential.
When $w \sim N/2$ as for random partial orders, it is exponential.

\begin{remark}[The Hauptvermutung]
The question of whether the causal order determines the manifold has been open since the founding of causal set theory~\cite{bombelli1987, sorkin2003, malament1977}.
Theorem~\ref{thm:polywidth} provides the first proved quantitative answer: a finite causal set is manifold-like if and only if $\gamma$ grows polynomially with the number of elements.
For $d$-dimensional spacetimes, $w \sim N^{(d-1)/d}$ gives $|\CC| = O(N^{2(d-1)})$\textemdash polynomial.
For random partial orders, $w \sim N/2$ gives $|\CC|$ exponential.
The gap is not a conjecture; it is a machine-verified theorem.
\end{remark}


%% ================================================================
\section{Non-Abelian Gauge Connection and Holonomy}\label{sec:holonomy}
%% ================================================================

The vanishing of abelian \v{C}ech $H^1$ for the Benincasa--Dowker curvature presheaf~\cite{benincasa2010} indicates that discrete curvature is fundamentally a connection phenomenon.
We construct a non-abelian gauge connection on $\CSpec$ whose holonomy detects curvature directly.

\subsection{Transition matrices}

Given a spatial-wedge cover $\{U_s\}$ of a discretized spacetime, the interval-count matrix $A_s$ with entries $(A_s)_{x,k} = N_k(x, U_s)$\textemdash counting $(k{+}2)$-element causal intervals above $x$ within $U_s$\textemdash provides a local frame for the causal geometry.
On overlaps $U_s \cap U_{s+1}$, the \emph{transition matrix}
\[
T_{s,\,s+1} = A_s^+\, A_{s+1}
\]
transforms the interval description from frame $s$ to frame $s+1$, where $A_s^+$ denotes the Moore--Penrose pseudoinverse restricted to the overlap.

\subsection{The Wilson loop}

\begin{definition}
For a closed loop $\gamma = (s_0, s_1, \ldots, s_n = s_0)$ in the cover, the \emph{holonomy} is $\Hol(\gamma) = \prod_i T_{s_i, s_{i+1}}$ and the \emph{Wilson loop} is
\[
W_\gamma = \frac{\operatorname{tr}(\Hol(\gamma))}{\dim(\Hol(\gamma))}.
\]
\end{definition}

The Wilson loop is gauge-invariant: it does not depend on the choice of starting wedge.

\begin{theorem}\label{thm:holonomy}
On discretized cylinders $S^1 \times \R$:

\medskip
\begin{center}
\begin{tabular}{cccc}
\toprule
Spatial circumference & Temporal depth & $W$ & $\operatorname{tr}(\Hol)$ \\
\midrule
4  & 3 & $2/5$ & 2 \\
5  & 3 & $2/5$ & 2 \\
6  & 4 & $3/5$ & 3 \\
8  & 4 & $3/5$ & 3 \\
10 & 3 & $2/5$ & 2 \\
\bottomrule
\end{tabular}
\end{center}
\medskip

\noindent The Wilson loop values are exact rationals, and
\[
\operatorname{tr}(\Hol) = n_{\mathrm{time}} - 1
\]
for every tested configuration.
This integer is independent of spatial circumference.
\end{theorem}

The independence from spatial circumference is physically correct: on a flat cylinder, the spatial circle carries no intrinsic curvature; the holonomy measures the parallel transport cost of the temporal extent.
The eigenvalue spectrum of $\Hol$ exhibits one eigenvalue near unity\textemdash the flat direction, corresponding to the translation-invariant count $N_0$ of direct causal links\textemdash with remaining eigenvalues decaying toward zero.


\subsection{Curvature discrimination}

The \v{C}ech $H^2$ of the Benincasa--Dowker curvature presheaf discriminates spatial topology where $H^1$ cannot:

\medskip
\begin{center}
\begin{tabular}{lccc}
\toprule
Configuration & $\dim H^0$ & $\dim H^1$ & $\dim H^2$ \\
\midrule
Cylinder $5 \times 3$ & 15 & 0 & 40 \\
Strip $5 \times 3$    & 15 & 0 & 20 \\[2pt]
Cylinder $6 \times 4$ & 24 & 0 & 144 \\
Strip $6 \times 4$    & 24 & 0 & 88 \\
\bottomrule
\end{tabular}
\end{center}
\medskip

\noindent The gap $\Delta H^2 = 20$ at scale $5 \times 3$ and $\Delta H^2 = 56$ at scale $6 \times 4$ grows with the size of the topological feature.


%% ================================================================
\section{The Arithmetic Bridge}\label{sec:arithmetic}
%% ================================================================

The divisibility lattice $(\Z^+, \mid)$ is a locally finite partial order and therefore a causal set in the formal sense.

\begin{theorem}[Arithmetic Bridge \lv{arithmetic\_bridge}\,$[\checkmark_L]$]\label{thm:bridge}
\begin{enumerate}[nosep]
\item $(\Z^+, \mid)$ is a causal algebra under the divisibility order.
\item $\mu * \zeta = 1$ in the arithmetic function ring.
\item $(\mu * \zeta)(n) = [n = 1]$.
\item $\mu(p) = -1$ for every prime $p$.
\item Every prime $p \mid n$ generates a causally prime ideal in $\CSpec$, recovering $\Spec(\Z)$.
\end{enumerate}
This is not an analogy.
The causal-algebraic framework contains classical number theory as a special case.
\end{theorem}

\begin{remark}[Holonomy on the divisibility lattice]
On the truncated lattice $(\Z^+ \cap [1, N], \mid)$, complex holonomy eigenvalues appear at $N \approx p^2$ for each prime $p$, with imaginary accumulation $A(p) \sim C \cdot p^{-\alpha}$, $\alpha \approx 3.27$ ($R^2 = 0.998$; computed on NVIDIA H100).
The partial sums $\sum A(p)$ converge to approximately $4.90$.
Four obstacles\textemdash the infinite limit, analytic continuation, the self-adjoint operator, and the missing Frobenius endomorphism\textemdash separate this observation from the Riemann Hypothesis.
\end{remark}


%% ================================================================
\section{Formal Verification}\label{sec:verification}
%% ================================================================

The framework is formalized in Lean~4 with Mathlib~v4.28.0: 18 files totaling 3{,}419 lines, with zero unproved assertions and zero custom axioms.
Every theorem depends only on \texttt{propext}, \texttt{Classical.choice}, and \texttt{Quot.sound}\textemdash the three standard kernel axioms.
The uniqueness theorem (Theorem~\ref{thm:uniqueness}) requires only \texttt{propext} and \texttt{Quot.sound}\textemdash it is constructive.
Computational results use \texttt{native\_decide} (\texttt{Lean.ofReduceBool}, \texttt{Lean.trustCompiler}), standard for verified computation.

\medskip
\begin{center}
\begin{tabular}{lrl}
\toprule
\textbf{Module} & \textbf{Lines} & \textbf{Key verified result} \\
\midrule
\texttt{CausalAlgebra}          & 221 & Causal ring, idempotents, corner algebras \\
\texttt{CausalPrimality}        & 226 & $\CSpec$, Recovery Theorem, NC collapse \\
\texttt{NoetherianRatio}        & 283 & $\gamma \geq 1$, total-order and antichain bounds \\
\texttt{ArithmeticBridge}       & 165 & $\mu * \zeta = 1$, primes $\to$ $\CSpec$ \\
\texttt{CSpecSheaf}             & 142 & Locality, convexity $\Rightarrow$ ring homomorphism \\
\texttt{OrderComplexCohomology} & 221 & $\delta^2 = 0$, $H^0$ characterization, cone kills $H^1$ \\
\texttt{Separation}             & 205 & $\gamma$ distinguishes Y-shape from Fork \\
\texttt{GridBound}              & 175 & $\gamma(\text{chain}) < \gamma(\text{grid}) < \gamma(\text{antichain})$ \\
\texttt{WidthBound}             & 166 & $\gamma \leq 2^w$ computational evidence \\
\texttt{WidthOneProof}          & 120 & $\gamma < 2$ for all total orders \\
\texttt{ChainRestriction}       & 133 & Convex $\to$ interval on chains \\
\texttt{Dilworth}               & 240 & Chain cover from inductive step \\
\texttt{Hauptvermutung}         & 158 & Formal conjecture statement \\
\texttt{Uniqueness}             & 149 & Causal primality is the unique maximal primality \\
\texttt{ConvexFactorization}    & $\sim$200 & Injective factorization, interval characterization \\
\texttt{BalancedBound}          & $\sim$150 & $f(m) \leq m^2+1$, balanced product bound \\
\texttt{PolynomialWidth}        & $\sim$200 & $|\CC(C)| = O(n^{2w})$ for fixed width \\
\texttt{InjectionBound}         & $\sim$120 & Injective map $\Rightarrow$ cardinality bound \\
\midrule
\textbf{Total}                  & \textbf{3{,}419} & \textbf{Zero unproved assertions, zero custom axioms} \\
\bottomrule
\end{tabular}
\end{center}

\medskip

The source code is available at \url{https://github.com/tomdif/causal-algebraic-geometry-lean}.

To the author's knowledge, this is the first instance of a new geometric framework being formally verified from its inception.


%% ================================================================
\section{Discussion and Open Problems}\label{sec:discussion}
%% ================================================================

\subsection{Correspondence with scheme theory}

The causal-algebraic framework provides partial orders the same structural depth that scheme theory provides commutative rings:

\medskip
\begin{center}
\begin{tabular}{lcl}
\toprule
Commutative Algebra & & Causal-Algebraic Geometry \\
\midrule
Commutative ring $R$           & $\longleftrightarrow$ & Causal algebra $(A, \{{\bf e}_\alpha\}, \leq)$ \\
$\Spec(R)$                     & $\longleftrightarrow$ & $\CSpec(A)$ \\
Localization $R_{\mathfrak{p}}$& $\longleftrightarrow$ & Corner algebra ${\bf e}_S A\, {\bf e}_S$ \\
Structure sheaf $\mathcal{O}$  & $\longleftrightarrow$ & Corner algebra sheaf \\
Sheaf cohomology               & $\longleftrightarrow$ & Order-complex cohomology \\
$\Spec(\Z)$                    & $\longleftrightarrow$ & $\CSpec(k[{\rm Div}(n)])$ \\
Noetherian condition            & $\longleftrightarrow$ & $\gamma \approx 1$ \\
Spectrum determined by ring     & $\longleftrightarrow$ & Spectrum determined by sheaf (Thm.~\ref{thm:uniqueness}) \\
\bottomrule
\end{tabular}
\end{center}


\subsection{Open problems}

\begin{enumerate}
\item \textbf{Sharpen the width bound.}
The polynomial width theorem gives $|\CC(C)| = O(n^{2w})$.
The exponential conjecture $\gamma \leq 2^w$ remains open and would give a tighter bound.
The chain restriction, decomposition injectivity, and Dilworth assembly are formally verified; what remains is a counting argument tighter than the balanced product bound.

\item \textbf{Continuum limit.}
Prove that $\CSpec$ and $\gamma$ converge as the poset is refined via Poisson sprinkling at increasing density, requiring a theory of pro-$\CSpec$ analogous to Grothendieck's pro-\'etale constructions.

\item \textbf{Causal Grothendieck site.}
Replace the Alexandrov topology with a Grothendieck topology whose covering sieves are families of causal embeddings, potentially yielding non-trivial abelian $H^1$.

\item \textbf{Higher-dimensional scaling.}
Compute $\gamma$ for Poisson sprinklings in $d$-dimensional Minkowski space for $d = 2, 3, 4$ and verify the predicted scaling.

\item \textbf{The arithmetic programme.}
Characterize the $p^2$ resonance in the divisibility lattice holonomy and determine whether the convergent accumulation $A(p) \sim C \cdot p^{-3}$ admits a number-theoretic explanation.
\end{enumerate}


\section*{Acknowledgments}

The formal verification uses Mathlib~v4.28.0~\cite{mathlib}.
Computational experiments were performed on NVIDIA H100 80GB HBM3 GPUs.


\begin{thebibliography}{99}

\bibitem{bombelli1987}
L.~Bombelli, J.~Lee, D.~Meyer, and R.D.~Sorkin,
Spacetime as a causal set,
\emph{Phys.\ Rev.\ Lett.} \textbf{59} (1987), no.~5, 521--524.

\bibitem{benincasa2010}
D.M.T.~Benincasa and F.~Dowker,
The scalar curvature of a causal set,
\emph{Phys.\ Rev.\ Lett.} \textbf{104} (2010), 181301.

\bibitem{rota1964}
G.-C.~Rota,
On the foundations of combinatorial theory~I: Theory of M\"obius functions,
\emph{Z.\ Wahrscheinlichkeitstheorie} \textbf{2} (1964), 340--368.

\bibitem{grothendieck1960}
A.~Grothendieck,
\'El\'ements de g\'eom\'etrie alg\'ebrique,
\emph{Publ.\ Math.\ IH\'ES} (1960--1967).

\bibitem{deligne1974}
P.~Deligne,
La conjecture de Weil, I,
\emph{Publ.\ Math.\ IH\'ES} \textbf{43} (1974), 273--307.

\bibitem{connes1994}
A.~Connes,
\emph{Noncommutative Geometry}, Academic Press, 1994.

\bibitem{voevodsky2003}
V.~Voevodsky,
Motivic cohomology with $\Z/2$-coefficients,
\emph{Publ.\ Math.\ IH\'ES} \textbf{98} (2003), 59--104.

\bibitem{sorkin2003}
R.D.~Sorkin,
Causal sets: Discrete gravity,
in \emph{Lectures on Quantum Gravity} (A.~Gomberoff and D.~Marolf, eds.), Springer, 2003.

\bibitem{malament1977}
D.~Malament,
The class of continuous timelike curves determines the topology of spacetime,
\emph{J.\ Math.\ Phys.} \textbf{18} (1977), 1399--1404.

\bibitem{myrheim1978}
J.~Myrheim,
Statistical geometry,
CERN preprint TH-2538, 1978.

\bibitem{dilworth1950}
R.P.~Dilworth,
A decomposition theorem for partially ordered sets,
\emph{Ann.\ of Math.\ (2)} \textbf{51} (1950), 161--166.

\bibitem{lubotzky1988}
A.~Lubotzky, R.~Phillips, and P.~Sarnak,
Ramanujan graphs,
\emph{Combinatorica} \textbf{8} (1988), no.~3, 261--277.

\bibitem{mathlib}
The Mathlib Community,
\emph{Mathlib: The Lean Mathematical Library},
\url{https://github.com/leanprover-community/mathlib4}, 2024.

\end{thebibliography}

\end{document}
